\chapter{Analyse et spécification des besoins}

\section{Introduction}

L'analyse et la spécification des besoins constituent la première étape essentielle du cycle
de vie d'un projet. Ce chapitre vise à définir clairement la vision du projet, à identifier les
besoins fonctionnels et non fonctionnels, ainsi que les principales fonctionnalités attendues.
Il s'agit également de repérer les différents acteurs qui interagiront avec le système, de
planifier les versions futures, de constituer le backlog des tâches, et enfin d'établir une
planification initiale des phases ainsi qu'un calendrier prévisionnel de réalisation du projet.

%% ----------------------------------------------------------------
\section{Étude des besoins}

Tout au long de cette section, nous allons identifier et spécifier les besoins à satisfaire, qui
représentent les fonctionnalités à mettre en œuvre dans notre application.

\subsection{Identification des acteurs}

Un acteur est une représentation idéalisée d'un rôle joué par une entité externe. Il peut
avoir le droit de consulter le système et éventuellement de le modifier. Pour notre application,
nous avons identifié les acteurs suivants :

\begin{itemize}
    \item \textbf{Client} : utilisateur anonyme accédant via QR code depuis son smartphone,
    sans inscription préalable.

    \item \textbf{Client Fidèle} : client possédant un compte avec programme de points de fidélité.

    \item \textbf{Staff} : employé du café gérant les commandes et les demandes d'assistance
    depuis le tableau de bord.

    \item \textbf{Admin} : administrateur disposant des droits complets sur le système, les
    utilisateurs et les statistiques.

    \item \textbf{Système RSS} : acteur système externe alimentant automatiquement le flux
    d'actualités en temps réel.

    \item \textbf{Système Chatbot IA} : acteur système autonome traitant les requêtes
    conversationnelles en français, anglais et arabe, basé sur une architecture RAG
    (\textit{Retrieval-Augmented Generation}).
\end{itemize}

%% ----------------------------------------------------------------
\subsection{Besoins fonctionnels}

Un besoin fonctionnel désigne une action spécifique que le système doit pouvoir réaliser.
Notre solution propose un ensemble de fonctionnalités adaptées aux besoins de chaque
type d'utilisateur.

\paragraph{Client :}
\begin{itemize}
    \item Scanner le QR code de sa table pour créer une session anonyme
    \item Consulter le menu des produits disponibles par catégories
    \item Passer une commande et suivre son statut en temps réel
    \item Appeler un serveur en cas de besoin
    \item Interagir avec le chatbot Luna pour obtenir des informations sur le menu,
          les promotions ou les ingrédients, en français, anglais ou arabe
    \item Jouer à un jeu (quiz, puzzle ou mots croisés) pendant l'attente
    \item Consulter les actualités en temps réel via flux RSS
\end{itemize}

\paragraph{Client Fidèle :}
\begin{itemize}
    \item Toutes les fonctionnalités du Client
    \item Créer un compte personnel avec nom et numéro de téléphone
    \item Accumuler des points de fidélité à chaque commande ou session de jeu
    \item Échanger ses points contre des récompenses
\end{itemize}

\paragraph{Staff :}
\begin{itemize}
    \item S'authentifier pour accéder au tableau de bord
    \item Visualiser et gérer les commandes en cours (accepter, préparer, livrer)
    \item Répondre aux demandes d'assistance des clients
\end{itemize}

\paragraph{Admin :}
\begin{itemize}
    \item S'authentifier pour accéder au tableau de bord
    \item Consulter les statistiques globales de ventes et d'activité
    \item Gérer les produits, catégories et utilisateurs (CRUD complet)
    \item Configurer les paramètres du système
\end{itemize}

\paragraph{Système RSS :}
\begin{itemize}
    \item Récupérer et mettre à jour automatiquement les actualités en temps réel
          depuis des flux RSS externes
\end{itemize}

\paragraph{Système Chatbot IA :}
\begin{itemize}
    \item Comprendre les questions en langage naturel (français, anglais, arabe)
    \item Générer un vecteur d'embedding de la requête via le modèle
          \textit{multilingual-e5-small} (exécuté localement)
    \item Récupérer les produits et FAQ pertinents depuis PostgreSQL par similarité cosinus
    \item Générer des réponses personnalisées via le modèle LLaMA~3.1 (API Groq)
    \item Détecter automatiquement la langue de l'utilisateur et adapter la réponse
          en conséquence
\end{itemize}

%% ----------------------------------------------------------------
\subsection{Besoins non fonctionnels}

Les besoins non fonctionnels décrivent toutes les contraintes techniques, ergonomiques et
esthétiques auxquelles le système est soumis. Le tableau~\ref{tab:non_fonctionnels} présente
les exigences retenues pour notre système.

\begin{table}[H]
    \centering
    \caption{Besoins non fonctionnels du système}
    \label{tab:non_fonctionnels}
    \begin{tabularx}{\textwidth}{|L{4cm}|X|}
        \hline
        \textbf{Critère} & \textbf{Exigence} \\
        \hline
        Intégrité &
        Garantir la cohérence et l'intégrité des données stockées dans PostgreSQL \\
        \hline
        Ergonomie &
        Interfaces conviviales, utilisation simple et intuitive pour tous les profils \\
        \hline
        Performance &
        Temps de réponse API inférieur à 500~ms pour les opérations courantes \\
        \hline
        Rapidité &
        Exécution des requêtes approchant le temps réel \\
        \hline
        Sécurité &
        Authentification JWT, validation des entrées, protection des routes par rôle \\
        \hline
        Intelligence conversationnelle &
        Réponses générées via LLaMA~3.1 (Groq), latence cible $<$~3~secondes \\
        \hline
        Multilinguisme &
        Détection automatique FR/EN/AR, adaptation des réponses à la langue détectée \\
        \hline
        Précision contextuelle &
        Réponses basées uniquement sur les données réelles du menu PostgreSQL \\
        \hline
        Scalabilité &
        Support de 50+ conversations simultanées sans dégradation de performance \\
        \hline
    \end{tabularx}
\end{table}

%% ================================================================
\section{Planification du projet}

\subsection{Backlog produit}

Chaque élément du backlog est caractérisé par un identifiant unique, la fonctionnalité
concernée, une user story décrivant le besoin de l'utilisateur, une priorité et une
complexité estimée.

\begin{table}[H]
\centering
\caption{Backlog produit global}
\label{tab:backlog}

\small
\renewcommand{\arraystretch}{1.2}
\setlength{\tabcolsep}{6pt}

\begin{tabular}{|c|p{4cm}|p{6.5cm}|c|c|}
\hline

\textbf{ID} &
\textbf{Fonctionnalité} &
\textbf{User Story} &
\textbf{Priorité} &
\textbf{Complexité} \\
\hline

1 &
Authentification &
En tant que Staff ou Admin, je veux m'authentifier pour accéder au tableau de bord selon mon rôle. &
Haute &
Faible \\
\hline

2 &
Session QR Code &
En tant que Client, je veux scanner le QR code de ma table pour démarrer une session anonyme sans inscription. &
Haute &
Moyenne \\
\hline

3 &
Consultation du menu &
En tant que Client, je veux consulter les produits organisés par catégories et rechercher un article. &
Haute &
Moyenne \\
\hline

4 &
Gestion du panier &
En tant que Client, je veux ajouter des produits au panier, voir le total et confirmer ma commande. &
Haute &
Élevée \\
\hline

5 &
Suivi de commande &
En tant que Client, je veux suivre l'état de ma commande en temps réel jusqu'à sa livraison. &
Haute &
Moyenne \\
\hline

6 &
Appel serveur et facture &
En tant que Client, je veux appeler un serveur et générer une facture imprimable. &
Moyenne &
Faible \\
\hline

7 &
Dashboard Admin/Staff &
En tant qu'Admin ou Staff, je veux gérer les commandes, produits et statistiques. &
Haute &
Élevée \\
\hline

8 &
Gestion des utilisateurs &
En tant qu'Admin, je veux gérer les comptes Staff. &
Haute &
Faible \\
\hline

9 &
Divertissement &
En tant que client je veux accéder à des jeux et citations pendant l'attente. &
Moyenne &
Moyenne \\
\hline

10 &
Actualités RSS &
En tant que client je veux consulter des actualités en temps réel via RSS. &
Moyenne &
Faible \\
\hline

11 &
Fidélité &
En tant que client je veux accumuler des points et les échanger contre des récompenses. &
Moyenne &
Élevée \\
\hline

12 &
Chatbot IA &
 En tant que client je veux poser des questions au chatbot (FR/EN/AR). &
Moyenne &
Élevée \\
\hline

\end{tabular}
\end{table}

\subsection{Découpage des sprints}

Le projet est organisé en deux releases, chacune composée de deux sprints, permettant une
livraison incrémentale des fonctionnalités. La figure~\ref{fig:decoupage_sprints} illustre cette répartition.

\begin{figure}[H]
\centering
\begin{tikzpicture}[
    release/.style={
        rectangle, rounded corners=4pt,
        fill=NavyBlue!80, text=white,
        font=\bfseries\large,
        minimum width=1.8cm, minimum height=5.2cm,
        align=center
    },
    sprint/.style={
        rectangle, rounded corners=4pt,
        fill=NavyBlue!60, text=white,
        font=\bfseries,
        minimum width=2.4cm, minimum height=2.2cm,
        align=center
    },
    features/.style={
        rectangle, rounded corners=4pt,
        fill=NavyBlue!45, text=white,
        font=\small\bfseries,
        minimum width=4.2cm, minimum height=2.2cm,
        align=left,
        text width=3.8cm,
        inner sep=0.35cm
    },
    outer/.style={
        rectangle, rounded corners=6pt,
        fill=NavyBlue!20,
        inner sep=0.3cm
    }
]

% ══════════════════════════════════════════════════════════
%  RELEASE 1
% ══════════════════════════════════════════════════════════
\node[release] (r1) at (0,0) {Release\\[0.1cm]1};

\node[sprint] (s1) at ([xshift=2.95cm, yshift=1.45cm] r1.center) {Sprint 1};
\node[features] (f1) at ([xshift=3.75cm] s1.center) {
    \textbullet\ Authentification JWT\\[0.2cm]
    \textbullet\ Gestion des rôles
};

\node[sprint] (s2) at ([xshift=2.95cm, yshift=-1.45cm] r1.center) {Sprint 2};
\node[features] (f2) at ([xshift=3.75cm] s2.center) {
    \textbullet\ Session QR code\\[0.2cm]
    \textbullet\ Gestion menu
};

\begin{scope}[on background layer]
    \node[outer, fit=(r1)(s1)(f1)(s2)(f2)] {};
\end{scope}

% ══════════════════════════════════════════════════════════
%  RELEASE 2
% ══════════════════════════════════════════════════════════
\node[release] (r2) at (0,-5.6) {Release\\[0.1cm]2};

\node[sprint] (s3) at ([xshift=2.95cm, yshift=1.45cm] r2.center) {Sprint 3};
\node[features] (f3) at ([xshift=3.75cm] s3.center) {
    \textbullet\ Dashboard Admin\\[0.2cm]
};

\node[sprint] (s4) at ([xshift=2.95cm, yshift=-1.45cm] r2.center) {Sprint 4};
\node[features] (f4) at ([xshift=3.75cm] s4.center) {
    \textbullet\ Divertissement\\[0.2cm]
    \textbullet\ Chatbot IA \& RSS
};

\begin{scope}[on background layer]
    \node[outer, fit=(r2)(s3)(f3)(s4)(f4)] {};
\end{scope}

\end{tikzpicture}
\caption{Découpage des sprints}
\label{fig:decoupage_sprints}
\end{figure}

\section{Conclusion}
L'analyse approfondie des besoins et la planification rigoureuse en sprints nous ont permis
d'établir une feuille de route claire pour le développement de Coffee Time. Ce chapitre a
défini le périmètre fonctionnel du projet, posant ainsi les bases nécessaires pour entamer
la phase de conception technique et de réalisation.
